\documentclass[man,12pt,floatsintext]{apa7}
\usepackage[american]{babel}
\usepackage{csquotes} 
\usepackage[style=apa,sortcites=true,sorting=nyt,backend=biber]{biblatex}
\DeclareLanguageMapping{american}{american-apa}

\usepackage[T1]{fontenc} 
\usepackage{mathptmx}

% 定义信息
\title{APA Format Starter Paper, Tailored to Lab Reports But Adaptable to Other Writing Assignments} 
\shorttitle{APA Starter} 
\author{Your Name Here}

\begin{document}

% --- 自定义标题页 ---
\begin{titlepage}
\makeatletter
    \centering
    \vspace*{2cm}
    {\large Tallinn University}\\
    \vspace{0.5cm}
    {\large School of Digital Technologies} \\
    \vspace{0.5cm}
    {\large DTLGM.DT Digital Learning Games}

    \vfill

    {\huge \textbf{ \@title }}\\
    \vspace{2cm}
    {\large Master's Thesis}

    \vfill
    \vfill

    \begin{flushright}
        {\large Author: \@author} \\
        \vspace{0.3cm}
        {\large Supervisor: Martin Sillaots}
    \end{flushright}

    \vfill

    {\large Tallinn 2026}
\makeatother
\end{titlepage}

% --- APA 7 标准摘要页 ---
\newpage
\section*{\centering Abstract} % APA 7 要求摘要标题居中且加粗（section* 在 man 模式下默认加粗）
\noindent % 摘要首行不缩进
Insert abstract text here. This follows the APA 7th edition manual guidelines for formatting an abstract page when not using the default maketitle command. 

\vspace{0.5in} 
\noindent \textit{Keywords:} APA style, demonstration % Keywords 标签需斜体，首行不缩进

% --- 正文 ---
\newpage
\section{Introduction}
Introductory text goes here.

\section{Method}
\subsection{Participants}
Participant information goes here.

\begin{table}[h]
    \caption{Sample words from this hypothetical experiment.}
    \centering
    \begin{tabular}{cc}
        \hline 
         First word & Second word \\
         \hline
         Yeet & Yoink \\
         Hot & Lit \\
         \hline
    \end{tabular}
    \label{tab:table_words}
\end{table}

\section{Results}
Statistical words.

\section{Discussion}
Discussion words.

% \printbibliography

\end{document}